
%%%%%%%%%%%%%%%%%%%%%%%%%%%%%%%%%%%%%%%%
%           main body
%%%%%%%%%%%%%%%%%%%%%%%%%%%%%%%%%%%%%%%%


%-----------------------------------
%	TITLE PAGE
%-----------------------------------

\title[Linear Algebra]{\LARGE \bfseries 第二部分\quad 线性代数} % The short title appears at the bottom of every slide, the full title is only on the title page
\author[Real Symmetric Matrix and Quadratic Form] % Your institution as it will appear on the bottom of every slide, may be shorthand to save space
{\Large \bfseries \S 6. 实对称阵与二次型}
% \author{Singular Value Decomposition} % Your name
% \date{\today} % Date, can be changed to a custom date
\date{}

%------------------------------------------------

\begin{document}


%------------------------------------------------

\begin{frame}

\titlepage % Print the title page as the first slide

\end{frame}

%------------------------------------------------

\begin{frame}
\frametitle{内容提要} % Table of contents slide, comment this block out to remove it
\tableofcontents % Throughout your presentation, if you choose to use \section{} and \subsection{} commands, these will automatically be printed on this slide as an overview of your presentation
\end{frame}

%-----------------------------------
%	PRESENTATION SLIDES
%-----------------------------------

%------------------------------------------------

\section{实对称矩阵的对角化}

%------------------------------------------------

\begin{frame}

\begin{block}{代数基本定理，Fundamental Theorem of Algebra}
任何一个复系数的单变量$n$次多项式，在复数集中都恰好有$n$个根（记重数）。
\end{block}

\vspace{2em}
\pause

对于一个$n$阶实方阵$A$来说，它的特征多项式$f(\lambda)$是一个$n$次实系数（从而也是复系数）的单变量多项式。因此有$n$个复根。

\end{frame}

%------------------------------------------------

\begin{frame}

\begin{block}{复矩阵}
一个复矩阵指的是一个矩阵$(a_{ij})$, 它的系数$a_{ij}$是复数。复矩阵也有加法、乘法以及复数的数乘。

\pause

\begin{itemize}
\item 复共轭转置：对于复矩阵$M = (a_{ij})$, 定义其共轭转置矩阵为$M^* = (a'_{ij})$, 其中$a'_{ij} = \overline{a}_{ji}$为$a_{ji}$的复共轭。
\end{itemize}
\end{block}

\vspace{1em}
\pause

\begin{block}{复向量的内积}
令$\alpha = (x_1,\ldots,x_n)^T,$ $\beta = (y_1,\ldots,y_n)^T$为两个$n$维复向量，则$\alpha$与$\beta$的内积被定义为
\[\langle \alpha, \beta \rangle = \alpha^*\beta = \overline{x}_1y_1 + \cdots + \overline{x}_ny_n.\]
于是，复向量$\alpha$范数（或长度）的定义为
\[\Vert \alpha \Vert = \sqrt{\langle \alpha, \alpha \rangle} = \sqrt{\overline{x}_1x_1 + \cdots + \overline{x}_nx_n}\]
\end{block}

\end{frame}

%------------------------------------------------

\begin{frame}

\begin{block}{实对称阵的特征值}
{\bfseries 命题}：$n$阶实对称方阵的有$n$个实的特征值。

\vspace{1em}
\pause

{\bfseries 证明}：把$n$阶实对称方阵$A$视作复系数的方阵，其特征多项式$f(\lambda)$视作一个复系数的单变量多项式，因此有$n$个复根。任取其中一个复根$\lambda$, 依定义存在复系数的$n$维非零向量$\alpha$，使得
\[A\alpha = \lambda\alpha.\]
于是
\[\alpha^*A\alpha = \lambda \alpha^*\alpha = \lambda \Vert \alpha \Vert^2.\]
又有
\[\alpha^*A\alpha = \alpha^*A^*\alpha = (A\alpha)^*\alpha = \overline{\lambda}\alpha^*\alpha = \overline{\lambda} \Vert \alpha \Vert^2.\]
因此$\lambda = \overline{\lambda}$,即$\lambda$为一个实数。
\end{block}

\end{frame}

%------------------------------------------------

\begin{frame}

\begin{block}{实对称阵的特征向量}
任取$n$阶实对称方阵$A$的一个特征值（实数），以及对应于它的一个特征向量$\alpha$. $\alpha$是如下的实系数齐次线性方程组的解
\[(\lambda I_n - A)\mathbf{x} = \mathbf{0}\]
于是$\alpha$是一个$n$维实向量。
\end{block}

\end{frame}

%------------------------------------------------

\begin{frame}

\begin{block}{实对称阵的特征向量}
之前已经证明了，属于不相等的特征值的特征向量是线性无关的。对实对称阵，有更强的结果：

\vspace{1em}
\pause

{\bfseries 命题}：实对称阵属于不相等的特征值的特征向量是正交的。

\vspace{1em}
\pause

{\bfseries 证明}：设$\alpha_1, \alpha_2$为分别属于实对称阵$A$的两个不相等特征值$\lambda_1, \lambda_2$, 即
\[A\alpha_1 = \lambda_1\alpha_1, \; A\alpha_2 = \lambda_2\alpha_2.\]
于是，我们有
\[\begin{cases}
\alpha_2^T A\alpha_1 = \alpha_2^T\lambda_1\alpha_1 = \lambda_1 \langle \alpha_2, \alpha_1 \rangle \\
\alpha_2^T A\alpha_1 = (\alpha_2^T A^T)\alpha_1 = (A\alpha_2)^T\alpha_1 = \lambda_2 \langle \alpha_2, \alpha_1 \rangle
\end{cases}\]
由于$\lambda_1 \neq \lambda_2$, 因此$\langle \alpha_2, \alpha_1 \rangle = 0.$
\end{block}

\end{frame}

%------------------------------------------------

\begin{frame}

\begin{block}{实对称阵的对角化}
如果$n$阶实对称阵的$n$个特征值互不相等，那么它可以对角化。并且还可以进一步要求选取出来的$n$个特征向量是标准正交的（Orthonormal），即向量间相互正交，而且每个向量长度都等于1.

\vspace{1em}
\pause

对于有重复特征值的情况，如果还是有$n$个线性无关的特征向量，则依然可以对角化。在每个特征子空间（即某个特征值对应的所有特征向量并上零向量）内，还可以用Gram-Schmidt正交化程序，得到标准正交基。

\vspace{1em}
\pause

有$n$个特征值的方阵相似于上三角方阵，而对称的上三角方阵只能是对角阵。
\end{block}

\end{frame}

%------------------------------------------------

\begin{frame}

\begin{block}{Gram-Schmidt正交化程序}
{\bfseries 目标}：从任意一组基$\alpha_1, \ldots, \alpha_n$出发构造一组标准正交基$\varepsilon_1, \ldots, \varepsilon_n$

\vspace{1em}
\pause

{\bfseries 方法}：归纳地构造。
\begin{enumerate}
\item 令$\varepsilon_1 = \alpha_1 / \Vert \alpha_1 \Vert$.
\item 假设已经有$k$个标准正交的向量组$\varepsilon_1, \ldots, \varepsilon_k$, 使得
\[\operatorname{span}\{ \varepsilon_1, \ldots, \varepsilon_k \} = \operatorname{span}\{ \alpha_1, \ldots, \alpha_k \},\]
令$\varepsilon' = \alpha_{k+1} + t_1\varepsilon_1 + \cdots + t_k\varepsilon_k,$ 从条件
\[0 = \langle \varepsilon', \varepsilon_i \rangle = \langle \alpha_{k+1}, \varepsilon_i \rangle + t_i\]
定出$t_i$, $i=1,\ldots,k$. 最后令
\[\varepsilon_{k+1} = \varepsilon' / \Vert \varepsilon' \Vert.\]
\end{enumerate}
\end{block}

\end{frame}

%------------------------------------------------

\begin{frame}

\begin{block}{实对称阵的对角化}

\vspace{1em}

{\bfseries 定理}：实对称阵正交相似于对角阵。

\vspace{1em}

也就是说，对于任意实对称阵$A$, 存在正交阵$\Omega$使得
\[\Omega^T A \Omega = \begin{pmatrix} \lambda_1 & & \\ & \ddots & \\ & & \lambda_n \end{pmatrix} = \Lambda\]

\begin{itemize}
\item 正交阵：即满足$\Omega^T\Omega = I_n$的方阵$\Omega$.
\end{itemize}
\end{block}

\end{frame}

%------------------------------------------------

\begin{frame}

\begin{block}{实对称阵正交对角化的例子}
我们来正交对角化实对称方阵$A = \begin{pmatrix} 1 & -2 & 0 \\ -2 & 2 & -2 \\ 0 & -2 & 3 \end{pmatrix}.$

\vspace{1em}
\pause

解：先求特征多项式
\begin{equation*}
\det(\lambda I_3 - A) = \det \begin{pmatrix} \lambda-1 & 2 & 0 \\ 2 & \lambda-2 & 2 \\ 0 & 2 & \lambda-3\end{pmatrix} =(\lambda+1)(\lambda-2)(\lambda-5)
\end{equation*}
因此$A$的特征值为$\lambda_1=-1,\lambda_2=2,\lambda_3=5$.

\begin{itemize}
\item 对于$\lambda_1$，由方程组$(-I_3-A) \mathbf{x} = \mathbf{0}$，解得其对应的单位特征向量为$\alpha_1 = (\frac{2}{3},\frac{2}{3},\frac{1}{3})^T$;
\item 对于$\lambda_2$，由方程组$(2I_3-A) \mathbf{x} = \mathbf{0}$, 解得其对应的单位特征向量为$\alpha_2 = (-\frac{2}{3}, \frac{1}{3}, \frac{2}{3})^T$;
\end{itemize}
\end{block}

\end{frame}

%------------------------------------------------

\begin{frame}

\begin{block}{实对称阵正交对角化的例子（续）}
\begin{itemize}
\item 对于$\lambda_3$，由方程组$(5I_3 - A) \mathbf{x} = \mathbf{0}$, 解得其对应的单位特征向量为$\alpha_3 = (\frac{1}{3},-\frac{2}{3},\frac{2}{3})^T$.
\end{itemize}

\vspace{2em}

于是，令
\[P = \begin{pmatrix} \frac{2}{3} & \frac{2}{3} & \frac{1}{3} \\ -\frac{2}{3} & \frac{1}{3} & \frac{2}{3} \\
\frac{1}{3} & -\frac{2}{3} & \frac{2}{3} \end{pmatrix},\]
那么
\[
P^TAP = \Lambda = \begin{pmatrix} -1 & 0 & 0 \\ 0 & 2 & 0 \\ 0 & 0 & 5 \end{pmatrix}.
\]
\end{block}

\end{frame}

%------------------------------------------------

\section{二次型}

%------------------------------------------------

\begin{frame}

\begin{block}{二次型的定义}
$\mathbb{R}$上的$n$个未定元$x_1,\ldots,x_n$的二次型为函数
\[Q(\mathbf{x}) = \mathbf{x}^TA\mathbf{x},\]
其中$A = (a_{ij})$为对称方阵，$\mathbf{x} = (x_1,\ldots,x_n)^T$为未定元的向量。

\pause

展开写即有
\begin{align*}
Q(\mathbf{x}) & = (x_1,\ldots,x_n) \begin{pmatrix} a_{11} & \cdots & a_{1n} \\ \vdots & & \vdots \\ a_{21} & \cdots & a_{nn} \end{pmatrix} \begin{pmatrix} x_1 \\ \vdots \\ x_n \end{pmatrix} \\
& = \sum\limits_{1\leqslant i,j \leqslant n} a_{ij}x_ix_j = \sum\limits_{i=1}^{n} a_{ii}x_i^2 + 2\sum\limits_{i\neq j} a_{ij}x_ix_j
\end{align*}

\pause

\begin{itemize}
\item 二次型 $\Longleftrightarrow$ $n$个未定元的二次齐次多项式 $\Longleftrightarrow$ 实对称阵
\end{itemize}
\end{block}

\end{frame}

%------------------------------------------------

\begin{frame}

\begin{block}{二次型的例子}
设$n$个变量的实值函数$f$在点$\mathbf{x} = (x_1,\ldots,x_n)$的一个邻域$U(\mathbf{x})$上有直到2阶的连续偏导数，且所有1阶偏导数在点$\mathbf{x}$的值都等于0。考虑$f$在$\mathbf{x}$处的泰勒展开式，如下
\[f(\mathbf{x} + \mathbf{h}) = f(\mathbf{x}) + \dfrac{1}{2!} \sum\limits_{i,j=1}^n \dfrac{\partial^2 f}{\partial x_i \partial x_j}(\mathbf{x})h_ih_j + r_3(\mathbf{x}; \mathbf{h})\]
其中的二次项
\[\sum\limits_{i,j=1}^n \dfrac{\partial^2 f}{\partial x_i \partial x_j}(\mathbf{x})h_ih_j = \mathbf{h}^T \underbrace{\left( \dfrac{\partial^2 f}{\partial x_i \partial x_j} \right)}_{H(f), \text{黑塞矩阵}} \mathbf{h}\]
就是一个关于增量向量$\mathbf{h}$的的二次型。
\end{block}

\end{frame}

%------------------------------------------------

\begin{frame}

\begin{block}{二次型的变量替换}
通过变量替换$\mathbf{y} = P^{-1}\mathbf{x}$, 有
\[Q(\mathbf{x}) = \mathbf{x}^T A \mathbf{x} = \mathbf{y}^T P^TAP \mathbf{y} = Q'(\mathbf{y})\]

\pause

通过这样的变量替换，能使二次型的形式尽量简单。例如：
\begin{displaymath}
Q(x_1,x_2) = ax_1^2 + 2bx_1x_2 + cx_2^2 = a\left( x_1 + \dfrac{b}{a}x_2 \right)^2 + \left( c - \dfrac{b^2}{a} \right) x_2^2
\end{displaymath}
亦即令\(\displaystyle \begin{pmatrix} y_1 \\ y_2 \end{pmatrix} = \begin{pmatrix} 1 & \frac{b}{a} \\ 0 & 1 \end{pmatrix} \begin{pmatrix} x_1 \\ x_2 \end{pmatrix}\), 有
\begin{align*}
Q(x_1,x_2) & = (y_1,y_2) \begin{pmatrix} 1 & 0 \\ -\frac{b}{a} & 1 \end{pmatrix} \begin{pmatrix} a & b \\ b & c \end{pmatrix} \begin{pmatrix} 1 & -\frac{b}{a} \\ 0 & 1 \end{pmatrix} \begin{pmatrix} y_1 \\ y_2 \end{pmatrix} \\
& = (y_1,y_2) {\color{red} \begin{pmatrix} a & 0 \\ 0 & c - \frac{b^2}{a} \end{pmatrix}} \begin{pmatrix} y_1 \\ y_2 \end{pmatrix}
\end{align*}
\end{block}

\end{frame}

%------------------------------------------------

\begin{frame}

\begin{block}{二次型变量替换的一般结果}
设$A$为$n$阶实对称方阵，那么存在可逆方阵$P$使得
\[P^T A P = \begin{pmatrix} a_1 & & \\ & \ddots & \\ & & a_n \end{pmatrix}\]
\end{block}

\pause

\begin{itemize}
\item $B = P^TAP$被称作$B$与$A$相合（Cogredient或Congruent）
\item 任意一个实对称阵都相合于形如以下的方阵
\[\begin{pmatrix} I_p & & \\ & -I_q & \\ & & 0 \end{pmatrix}\]
自然数$p,q$由$A$唯一决定，分别被称作其正惯性指数，负惯性指数。
\end{itemize}

\end{frame}

%------------------------------------------------

\begin{frame}

\begin{block}{实对称阵的相似与相合}
\begin{itemize}
\item 一般来说，对称阵才谈论相合；而任何方阵都可以谈论相似
\item 可对角化方阵对角化得到的对角方阵，在不计对角线元素顺序的情况下，是唯一的；实对称阵则能通过多种相合变换，得到不同的对角阵。
\item 相似与相合都来自坐标变换（或者说变量替换），前者来自于化简线性变换的矩阵表达，后者来自于化简二次型。
\item 当要求坐标变换的过渡矩阵$P$是正交阵的时候
\[P^T A P = P^{-1}AP\]
\end{itemize}
\end{block}

\end{frame}

%------------------------------------------------

\section{正定实对称阵}

%------------------------------------------------

\begin{frame}

\begin{block}{正定实对称阵的定义}
设$Q(\mathbf{x}) = \mathbf{x}^T A \mathbf{x}$为一个二次型，$A$为实对称阵。如果对于任意的非零向量$\mathbf{x}$都有
\[\mathbf{x}^T A \mathbf{x} > 0\]
则称$A$为一个正定（Positive Definite）实对称阵。
\end{block}

\vspace{1em}
\pause

\begin{block}{正定实对称阵的判别方法}
\begin{itemize}
\item $A$相合于单位阵，即存在可逆方阵$P$, 使得$A = P^TP$.
\item $A$的顺序主子式，即$\det (a_{ij})_{k\times k}$, $k = 1,2,\ldots,n$, 都大于0.
\item $A$的所有主子式，即$A$的第$i_1,\ldots,i_k$行与第$i_1,\ldots,i_k$列交点处元素按原来顺序排列得到的方阵的行列式，都大于0.
\end{itemize}
\end{block}

\end{frame}

%------------------------------------------------

\begin{frame}

\begin{block}{其他类型的实对称阵}
\begin{itemize}
\item 负定（Negative Definite）实对称阵：$\mathbf{x}^T A \mathbf{x} < 0$对所有非零的$\mathbf{x}$成立。\newline
\item 半正定（Semi-Positive）实对称阵：$\mathbf{x}^T A \mathbf{x} \geqslant 0$对所有$\mathbf{x}$成立。\newline
\item 半负定（Semi-Negative）实对称阵：$\mathbf{x}^T A \mathbf{x} \leqslant 0$对所有$\mathbf{x}$成立。\newline
\item 不定（Indefinite）实对称阵：$\mathbf{x}^T A \mathbf{x}$对有的$\mathbf{x}$取正值，对有的$\mathbf{x}$取负值。
\end{itemize}
\end{block}

\end{frame}

%------------------------------------------------

\begin{frame}

\begin{block}{在微积分中的应用}
设$n$个变量的实值函数$f$在点$\mathbf{x} = (x_1,\ldots,x_n)$的一个邻域$U(\mathbf{x})$上有直到2阶的连续偏导数，且所有1阶偏导数在点$\mathbf{x}$的值都等于0。考虑$f$在$\mathbf{x}$处的黑塞矩阵
\[\left( \dfrac{\partial^2 f}{\partial x_i \partial x_j} \right)_{n\times n}\]
\begin{itemize}
\item 正定时，$f$在点$\mathbf{x}$取得严格局部极小值。
\item 负定时，$f$在点$\mathbf{x}$取得严格局部极大值。
\item 不定时，$f$在点$\mathbf{x}$没有极值。
\pause
\item 半（正或负）定时，$f$在点$\mathbf{x}$处极值情况比较复杂。
\end{itemize}
\end{block}

\end{frame}

%------------------------------------------------

% \begin{frame}

% \begin{block}{Vandermonde行列式的计算}
% \begin{enumerate}
% % \addtocounter{enumi}{0}
% \item 把第$n-1$行乘以$-a_1$加到第$n$行，再把第$n-2$行乘以$-a_1$加到第$n-1$行。按照这种顺序依次向上，直到用第$2$行乘以$-a_1$加到第$1$行。由行列式性质，$V_n$的值不变。
% \[
% V_n = \det \begin{pmatrix}
% 1 & 1 & \cdots & 1 \\
% 0 & a_2-a_1 & \cdots & a_n-a_1 \\
% \vdots & \vdots & & \vdots \\
% 0 & a_2^{n-2}(a_2-a_1) & \cdots & a_n^{n-2}(a_n-a_1) \\
% \end{pmatrix}
% \]
% \end{enumerate}
% \end{block}

% \end{frame}

%------------------------------------------------
% % closing page
% \begin{frame}
% \HUGE{\centerline{\bfseries {\color{gray} The} {\color{zkblue} End}}}

% \pause

% \vspace{1.2em}

% \Large{\centerline{\bfseries {\color{gray}Thank} \quad {\color{zkblue} You}}}

% \end{frame}

%----------------------------------------------------------------------------------------

\end{document}
